\chapter*{General Conclusion}
\addcontentsline{toc}{chapter}{General Conclusion}
\markboth{General Conclusion}{}
\label{ch:conclusion}

This thesis addressed the automated grading of diabetic retinopathy in a context where
screening needs grow much faster than the number of human graders, and where an automated
system is only acceptable if it is both accurate and safe for patients. Starting from a
review of a decade of deep-learning methods, we showed that every published system shares
three architectural limitations: it grades each image in isolation, it ignores the global
topology of the vascular network, and it must answer even when the image is ambiguous.
The whole work was organised around resolving these three limitations together.

The study began with the clinical context of the disease and the methodological tools
needed to address the three limitations: graph neural networks, persistent homology,
ordinal regression and calibrated abstention. We then validated the graph template outside
medicine, on microservice and network-security graphs, and gathered the first evidence,
with three independent methods, that structural and topological features of the retina
carry information that deep visual features miss. These results led to TOPORET, an
inductive population graph whose edges combine visual and topological similarity. Its
pre-specified tests confirmed that the graph brings a significant gain over independent
classification ($t_4=3.82$, $p=0.009$) and that topological features are strongly
associated with severity ($F(4,4995)=1308.9$, $\eta^2=0.51$). TOPORET, however, remained
above the safety threshold. The final system, DOTS, added a retinal foundation model, an
ordinal head with guaranteed rank consistency and entropy-based abstention. On five public
benchmarks it meets the seven predefined criteria at the point estimate (QWK $=0.951$,
Sev-NR $=1.9\%$, $0\%$ rank violations), and an exhaustive enumeration showed that none of
its six reduced versions does. Finally, the system was engineered as medical-device
software whose design maps onto the IEC~62304 and ISO~14971 standards and the EU AI Act.

The three research questions of the thesis can therefore be answered as follows. Does
the relational context between patients carry information that a per-image classifier
cannot recover? Yes, provided that the graph is built with a similarity that includes the
topology of the vessels; a visual similarity alone is not enough. Is the global topology
of the vascular network related to the severity of the disease? Yes, strongly, with the
largest effects for the loops that appear with neovascularisation. Can a single
architecture meet all the predefined criteria of accuracy, calibration and safety, and are
all of its components needed? Yes at the point estimate, and none of its reduced versions
does, although the confidence interval of the safety criterion still calls for prospective
confirmation.

The contributions of this research are both scientific and practical. On the scientific
side, the thesis shows that the three limitations of automated DR grading can be resolved
in a single architecture, and that the three components are complementary rather than
interchangeable: within the evaluated family, removing any of them loses the safety
criterion. It also introduces the use of persistent-homology descriptors of the vessel
skeleton inside the similarity of a population graph, and a validation discipline in which
every hypothesis, test and threshold is fixed before the experiment. On the practical side,
DOTS grades an image in $31.4$~ms on a single GPU, defers only $6.2\%$ of images to human
graders, explains each decision with the similar cases and topological features that
support it, and is designed to be verified and certified. Under explicit prevalence
assumptions, it would avoid about $600$ missed sight-threatening cases per year in a
programme of $500{,}000$ patients. A CPU model developed along the way (SAMF-GB) also
offers a low-cost alternative for settings without GPUs.

Beyond its results, the research programme taught several methodological lessons that may
be useful to other work on medical AI. First, accuracy and safety must be measured
separately: every system in this thesis gained QWK long before it crossed the safety
threshold, and optimising the former does not guarantee the latter. Second, validating an
architectural idea outside medicine, where data are plentiful and labels unambiguous, made
the medical experiments easier to design and to interpret. Third, fixing hypotheses, tests
and thresholds before running the experiments turned some results that could have been
discarded, such as the non-significant visual-only graph, into informative findings.
Fourth, designing for regulation from the start cost little and avoided late changes to the
architecture. Finally, the most useful component for safety was not the most accurate one:
entropy-based abstention adds less than one point of QWK, but it is the step that brings
the non-referral rate below the threshold.

\phantomsection\label{sec:limits-gc}

These results must be read within clear limits. All experiments are retrospective, on
curated public datasets; the upper bound of the 95\% confidence interval of Sev-NR
($2.4\%$) is still above the $2\%$ threshold, so statistical compliance with the safety
criterion is not yet established. Performance is lower on the smallest dataset (IDRiD,
$516$ images), where the population graph is sparse. None of the datasets provides
demographic information, so fairness across age, sex and ethnicity could not be assessed.
The population graph uses a single visit per patient and ignores the progression of the
disease over time. Finally, the computation of topological features is the main latency
bottleneck, and the integration with hospital information systems (FHIR adapter) is
designed but not yet implemented.

These limits open several directions for future work. The most immediate is a
prospective, multi-site reader study, pre-registered, to confirm the safety of DOTS in
real screening conditions and to reduce the uncertainty on Sev-NR; it is the first step of
the four-phase translation roadmap of \cref{sec:roadmap}, which leads to a randomised
trial and to CE marking and FDA submission. On the methodological side, learnable
topological layers could replace the hand-crafted persistence vector and possibly improve
accuracy, provided that they remain interpretable. A temporal population graph, linking the
successive visits of each patient, could exploit the speed of progression at the difficult
Grade~2/3 boundary. Federated construction of the population graph would allow several
hospitals to collaborate without sharing patient data. Smartphone-based fundus capture,
connected to the DOTS inference service, could extend screening to low-resource settings.
Finally, an explanation layer based on large language models could turn the logged evidence
into narratives adapted to clinicians and patients, under strict safeguards that prevent
it from contradicting the grade or the decision to abstain.

More broadly, the three-axis framework -- relational graph, structural topology, and
ordinal prediction with safe abstention -- is not specific to the retina. It can serve as a
template for any graded disease with an ordinal severity scale, a structural anatomical
prior and a safety threshold, such as age-related macular degeneration or glaucoma.

In summary, this work contributes to making automated screening of diabetic retinopathy
not only more accurate, but also safer, more transparent and closer to clinical use. The
code, models and configurations are released for independent reproduction, and the
remaining steps towards deployment are clearly identified. The prevention of blindness at
the scale of a population will depend on the ability to carry such systems from
computational validation to everyday clinical practice.
